\scanpage{150}
% archive sheet number: 149
développements qui suivent ne font intervenir que les $\unsure{\mathcal V|\mathscr V}_\ell(S)$, mais en pratique on partira souvent avec des éléments de $\unsure{\mathcal M|\mathscr M}_\ell(S)$ [ou de façon équivalente par dualité de Pontrjagin, de $\unsure{\mathcal I|\mathrm I}_\ell(S)$].

Soit $S'\longrightarrow S$, il s'en déduit un \inserted{\unsure{diagramme de}} foncteurs \unsure{essentiellement|partiellement} commutatif
\diagram{\begin{tikzcd}[ampersand replacement=\&]
\mathcal I_\ell(S)^\circ\arrow[r,"\sim"]\arrow[d] \& \mathcal M_\ell(S)\arrow[r]\arrow[d] \& \mathcal V_\ell(S)\arrow[d]\\
\mathcal I_\ell(S')^\circ\arrow[r,"\sim"] \& \mathcal M_\ell(S')\arrow[r] \& \mathcal V_\ell(S')
\end{tikzcd}}
les flèches verticales étant des flèches de changement de base. Quand $S$ et $S'$ sont connexes, $\xi'$ un pt géométrique de $S'$, $\xi$ son image dans $S$, d'où
\[
\pi'=\pi_1(S',\xi')\unsure{\text{---}|\longrightarrow}\pi=\pi_1(S,\xi)
\]
alors les foncteurs verticaux s'interprètent [en termes de \deleted{\unsure{$\pi$ et $\pi$}} Modules sur $\pi$ et $\pi'$] par restriction des scalaires

